\documentclass[10pt,a4paper]{amsart}
\usepackage[margin=19mm]{geometry}
\usepackage{amsmath,amssymb,mathtools}
\usepackage[scr=rsfs,cal=euler]{mathalfa}
\usepackage{xfrac}
\usepackage[unicode,hidelinks,bookmarks=false]{hyperref}
\newcommand{\ie}{i.e.\ }
\newcommand{\cf}{cf.\ }
\newcommand{\resp}{respectively\ }
\newcommand{\Subsection}{\subsection}
\providecommand{\nobreakdash}{\nobreak\hbox{-}}
\setlength{\emergencystretch}{2em}
\multlinegap0pt
\usepackage{tikz}
\usepackage{float}
\usepackage{mathtools}
\usepackage{amssymb}
\usepackage{dsfont}
\usepackage{xcolor}
\newtheorem{lem}{Lemma}
\DeclareMathOperator{\epi}{epi}
\title{On the sequential convergence of Lloyd's algorithms}
\author{L\'eo Portales, Elsa Cazelles, Edouard Pauwels}
\date{Source: arXiv:2405.20744v5 (CC BY 4.0)}
\begin{document}
\maketitle

\noindent\emph{Source and licence.} On the sequential convergence of Lloyd's algorithms. Version arXiv:2405.20744v5 (CC BY 4.0), distributed by arXiv under the Creative Commons Attribution 4.0 International licence, http://creativecommons.org/licenses/by/4.0/, as recorded in the arXiv licence metadata for this exact version and shown on the abstract page https://arxiv.org/abs/2405.20744v5. Full licence notice: \texttt{licenses/arxiv-2405.20744-CC-BY-4.0.txt}.\par\smallskip

\noindent\textit{Editorial scope.} Counted unit: the article's lem minmax (source0.tex lines 586--591). The article's complete original proof of it is delivered with it (source0.tex lines 937--957, one \texttt{proof} environment of 21 lines, which the article places away from the statement). No mathematics is written, repaired or re-derived: the statement and the proof are the article's own lines, in the article's own notation, and no retained line is substituted or re-flowed.  No further source-local context is needed: every cross-reference of every retained line resolves inside the two delivered spans.  Nothing was omitted from any retained span, so the package carries no omission record.  No retained line contains a citation, so the wrapper carries no bibliography and no bibliography entry.  No retained line contains a figure environment, an image-inclusion command or drawing code, so no image is delivered and the package declares no asset dependency.  Editorial formatting only, in addition: an AMS \texttt{amsart} preamble that restates the article's own declarations and macros used by the retained lines, namely the environments \texttt{lem} on separate counters, together with the standard packages those lines need.  The retained environments print in this document as Lemma 1 in order of appearance, whereas the article prints the same items with the numbering of its own full document.  The complete native source of the article as distributed in the arXiv e-print for this exact version is bundled unchanged as \texttt{sources/160061/source0.tex}.
\begin{lem}\label{minmax} \hfill
\begin{enumerate}
    \item[(i)] Let $(f_{j})_{j\in J}$ be a finite collection of definable functions. Then $\underset{j\in J}{\min}\: f_{j}$ and $\underset{j\in J}{\max}\: f_{j}$ are definable functions. 
    \item[(ii)] Let $f:\mathbb{R}^{d}\times \mathbb{R}^{n}\longrightarrow \mathbb{R}$ be a definable function and $E$ be a definable subset of $\mathbb{R}^{n}$. Then, provided they are attained, $\underset{w\in E}{\min}\: f(\cdot,w)$ and $\underset{w\in E}{\max}\: f(\cdot,w)$ are definable functions.
\end{enumerate}
\end{lem}

\begin{proof}[Proof of Lemma \ref{minmax}]

We prove the results (i) and (ii) for the minimum. The maximum case follows a symmetric proof using the hypograph.

First, we recall that by definition of the epigraph, $\epi \left( \underset{j\in J }{\min} \:f_{j}\right)=\bigcup_{j\in J}\epi(f_{j})$. We thus conclude for (i) since definable subsets are stable by finite unions.

For the second property (ii), we denote by $p:\mathbb{R}^{d+n+1}\longrightarrow \mathbb{R}^{d+1}$ the canonical projection on the first $d$ coordinates and the last coordinate. We aim at showing the following equality
\begin{equation*}
\epi\left(\underset{w\in E}{\min}\:f(.,w)\right)=p\left(\epi(f)\right).
\end{equation*}
Since the minimum is attained by hypothesis, we have
\begin{align*}
(x,y)\in \epi\left(\underset{w\in E}{\min}\: f(.\:,w)\right)
\quad\iff \quad & \underset{w\in E}{\min}\:f(x,w)\leq y\\
\iff \quad & \exists \ w'\in E\ \text{such that}\ f(x,w')\leq y\\
\iff \quad & (x,w',y)\in \epi(f)\\
\iff \quad & (x,y)\in \ p(\epi(f))\\
\end{align*}
which concludes the proof.

\end{proof}

\end{document}
